\section{Entry ranks in finite refinements}

\import{set/finite-interleave.tex}

The entry rank records where a point first enters a finite sequence of sets.
We use the last index as a default when the point has not entered earlier.
Thus the rank always exists, but a rank equal to the last index does not
distinguish membership in the last set from membership in none of the sets.
Under refinement, a positive entry rank becomes either an inserted index or
the following old index; rank zero stays zero. This alternative will control
the step values.

\begin{definition}\label{entry_rank_candidates}
    $\entryRankCandidates{c}{r}{x}
        =\{i\in r\mid x\in c(i)\}\union\{r\}$.
\end{definition}

\begin{definition}\label{entry_rank}
    $\entryRank{c}{r}{x}
        =\inters{\entryRankCandidates{c}{r}{x}}$.
\end{definition}

\begin{proposition}\label{entry_rank_candidates_iff}
    $i\in\entryRankCandidates{c}{r}{x}$ iff
        $i\in r$ and $x\in c(i)$ or $i=r$.
\end{proposition}
\begin{proof}
    Follows by \auto{entry_rank_candidates,union_iff,singleton_iff}.
\end{proof}

\begin{proposition}\label{entry_rank_candidates_inhabited}
    $\entryRankCandidates{c}{r}{x}$ is inhabited.
\end{proposition}
\begin{proof}
    $r\in\entryRankCandidates{c}{r}{x}$
        by \auto{entry_rank_candidates_iff}.
\end{proof}

\begin{proposition}\label{entry_rank_candidates_are_ordinals}
    Let $r$ be a natural number.
    Then every element of $\entryRankCandidates{c}{r}{x}$ is an ordinal.
\end{proposition}
\begin{proof}
    Follows by
        \auto{entry_rank_candidates_iff,natural_number_is_ordinal,elem_of_ordinal_is_ordinal}.
\end{proof}

\begin{proposition}\label{entry_rank_is_ordinal}
    Let $r$ be a natural number.
    Then $\entryRank{c}{r}{x}$ is an ordinal.
\end{proposition}
\begin{proof}
    Follows by
        \auto{entry_rank,inters_of_ordinals_is_ordinal,entry_rank_candidates_inhabited,entry_rank_candidates_are_ordinals}.
\end{proof}

\begin{proposition}\label{entry_rank_is_candidate}
    Let $r$ be a natural number.
    Then $\entryRank{c}{r}{x}\in\entryRankCandidates{c}{r}{x}$.
\end{proposition}
\begin{proof}
    Follows by
        \auto{entry_rank,inters_of_ordinals_elem,entry_rank_candidates_inhabited,entry_rank_candidates_are_ordinals}.
\end{proof}

\begin{proposition}\label{entry_rank_subseteq}
    Let $r$ be a natural number.
    Let $i\in\entryRankCandidates{c}{r}{x}$.
    Then $\entryRank{c}{r}{x}\subseteq i$.
\end{proposition}
\begin{proof}
    Follows by
        \auto{entry_rank,inters_of_ordinals_subseteq,entry_rank_candidates_inhabited,entry_rank_candidates_are_ordinals}.
\end{proof}

\begin{proposition}\label{entry_rank_subseteq_sentinel}
    Let $r$ be a natural number.
    Then $\entryRank{c}{r}{x}\subseteq r$.
\end{proposition}
\begin{proof}
    $r\in\entryRankCandidates{c}{r}{x}$
        by \auto{entry_rank_candidates_iff}.
    Follows by \obv{entry_rank_subseteq}.
\end{proof}

\begin{proposition}\label{entry_rank_is_natural}
    Let $r$ be a natural number.
    Then $\entryRank{c}{r}{x}$ is a natural number.
\end{proposition}
\begin{proof}
    $\entryRank{c}{r}{x}\in\entryRankCandidates{c}{r}{x}$
        by \obv{entry_rank_is_candidate}.
    Hence $\entryRank{c}{r}{x}\in r$ or
        $\entryRank{c}{r}{x}=r$
        by \auto{entry_rank_candidates_iff}.
    Follows by
        \obv{naturals_is_transitive_set,transitiveset_iff_subseteq,subseteq}.
\end{proof}

\begin{proposition}\label{entry_rank_in_successor}
    Let $r$ be a natural number.
    Then $\entryRank{c}{r}{x}\in\suc{r}$.
\end{proposition}
\begin{proof}
    $\entryRank{c}{r}{x}\in r$ or
        $\entryRank{c}{r}{x}=r$
        by \auto{entry_rank_is_candidate,entry_rank_candidates_iff}.
    Follows by \auto{suc_iff}.
\end{proof}

\begin{proposition}\label{entry_rank_no_earlier_entry}
    Let $r$ be a natural number.
    Let $i\in\entryRank{c}{r}{x}$.
    Then $x\notin c(i)$.
\end{proposition}
\begin{proof}
    Suppose $x\in c(i)$.
    $\entryRank{c}{r}{x}\subseteq r$
        by \obv{entry_rank_subseteq_sentinel}.
    Hence $i\in r$ by \obv{subseteq}.
    Thus $i\in\entryRankCandidates{c}{r}{x}$
        by \auto{entry_rank_candidates_iff}.
    Hence $\entryRank{c}{r}{x}\subseteq i$
        by \obv{entry_rank_subseteq}.
    $i$ and $\entryRank{c}{r}{x}$ are ordinals
        by \auto{entry_rank_is_ordinal,elem_of_ordinal_is_ordinal}.
    $i\subseteq\entryRank{c}{r}{x}$
        by \obv{ordinal_elem_implies_subseteq}.
    Hence $i=\entryRank{c}{r}{x}$
        by \obv{subseteq_antisymmetric}.
    Contradiction by \obv{in_irrefl}.
\end{proof}

\begin{proposition}\label{entry_rank_eq_from_first_candidate}
    Let $r$ be a natural number.
    Let $q\in\suc{r}$.
    Suppose $q=r$ or $q\in r$ and $x\in c(q)$.
    Suppose for every $i\in q$ we have $x\notin c(i)$.
    Then $\entryRank{c}{r}{x}=q$.
\end{proposition}
\begin{proof}
    $q\in\entryRankCandidates{c}{r}{x}$
        by \auto{entry_rank_candidates_iff}.
    Hence $\entryRank{c}{r}{x}\subseteq q$
        by \obv{entry_rank_subseteq}.
    Suppose not.
    $q$ and $\entryRank{c}{r}{x}$ are ordinals
        by \obv{natural_number_is_ordinal,entry_rank_is_ordinal,elem_of_ordinal_is_ordinal}.
    Hence $\entryRank{c}{r}{x}\in q$
        by \obv{ordinal_subseteq_implies_elem_or_eq}.
    $\entryRank{c}{r}{x}\in\entryRankCandidates{c}{r}{x}$
        by \obv{entry_rank_is_candidate}.
    $q\subseteq r$ by
        \obv{suc_iff,natural_number_is_ordinal,ordinal_elem_implies_subseteq,subseteq_refl}.
    Hence $\entryRank{c}{r}{x}\in r$ by \obv{subseteq}.
    Hence $x\in c(\entryRank{c}{r}{x})$
        by \auto{entry_rank_candidates_iff,in_irrefl}.
    Contradiction.
\end{proof}

\begin{proposition}\label{entry_rank_eq_zero_iff}
    Let $r$ be a natural number.
    Then $\entryRank{c}{r}{x}=\emptyset$ iff
        $r=\emptyset$ or $\emptyset\in r$ and $x\in c(\emptyset)$.
\end{proposition}
\begin{proof}
    If $\entryRank{c}{r}{x}=\emptyset$, then
        $r=\emptyset$ or $\emptyset\in r$ and $x\in c(\emptyset)$
        by \auto{entry_rank_is_candidate,entry_rank_candidates_iff}.
    Show that if $r=\emptyset$ or
        $\emptyset\in r$ and $x\in c(\emptyset)$, then
        $\entryRank{c}{r}{x}=\emptyset$.
    \begin{subproof}
        Assume $r=\emptyset$ or
            $\emptyset\in r$ and $x\in c(\emptyset)$.
        Then $\emptyset\in\entryRankCandidates{c}{r}{x}$
            by \auto{entry_rank_candidates_iff}.
        Hence $\entryRank{c}{r}{x}\subseteq\emptyset$
            by \obv{entry_rank_subseteq}.
        Follows by \obv{subseteq_emptyset_iff}.
    \end{subproof}
    Follows by \auto.
\end{proof}

\begin{proposition}\label{entry_rank_successor_predecessor}
    Let $r$ be a natural number.
    Suppose $\entryRank{c}{r}{x}=\suc{j}$.
    Then $j\in r$.
\end{proposition}
\begin{proof}
    $\entryRank{c}{r}{x}\subseteq r$
        by \obv{entry_rank_subseteq_sentinel}.
    $j\in\suc{j}$ by \obv{suc_intro_self}.
    Follows by \obv{subseteq}.
\end{proof}

\begin{proposition}\label{entry_rank_successor_candidate}
    Let $r$ be a natural number.
    Suppose $\entryRank{c}{r}{x}=\suc{j}$.
    Then $\suc{j}=r$ or
        $\suc{j}\in r$ and $x\in c(\suc{j})$.
\end{proposition}
\begin{proof}
    $\entryRank{c}{r}{x}\in\entryRankCandidates{c}{r}{x}$
        by \obv{entry_rank_is_candidate}.
    Follows by \auto{entry_rank_candidates_iff}.
\end{proof}

\begin{proposition}\label{entry_rank_eq_sentinel_outside_last}
    Let $r$ be a natural number.
    Suppose $c\in\funs{\suc{r}}{A}$.
    Suppose for every $i\in r$ we have $c(i)\subseteq c(r)$.
    Suppose $x\notin c(r)$.
    Then $\entryRank{c}{r}{x}=r$.
\end{proposition}
\begin{proof}
    $\entryRank{c}{r}{x}\in\entryRankCandidates{c}{r}{x}$
        by \obv{entry_rank_is_candidate}.
    Hence $\entryRank{c}{r}{x}\in r$ and
        $x\in c(\entryRank{c}{r}{x})$ or
        $\entryRank{c}{r}{x}=r$
        by \obv{entry_rank_candidates_iff}.
    Suppose not.
    Hence $\entryRank{c}{r}{x}\in r$ and
        $x\in c(\entryRank{c}{r}{x})$.
    Thus $c(\entryRank{c}{r}{x})\subseteq c(r)$.
    Hence $x\in c(r)$ by \obv{subseteq}.
    Contradiction.
\end{proof}

\subsection{Entry ranks under refinement}

\begin{proposition}\label{entry_rank_interleave_zero}
    Let $r$ be a natural number.
    Suppose $c\in\funs{\suc{r}}{A}$ and $m\in\funs{r}{A}$.
    Suppose $\entryRank{c}{r}{x}=\emptyset$.
    Then
        $\entryRank{\finiteInterleave{c}{m}}{r+r}{x}=\emptyset$.
\end{proposition}
\begin{proof}
    Let $d=\finiteInterleave{c}{m}$.
    Let $R=r+r$.
    $\dom{c}=\suc{r}$ and $\dom{m}=r$ by \obv{funs_elim}.
    $R$ is a natural number by \auto{num_addition_on_naturals}.
    We have $r=\emptyset$ or
        $\emptyset\in r$ and $x\in c(\emptyset)$
        by \obv{entry_rank_eq_zero_iff}.
    \begin{byCase}
    \caseOf{$r=\emptyset$.}
        Then $R=\emptyset$ by \auto{num_addition_zero_right}.
        Follows by \auto{entry_rank_eq_zero_iff}.
    \caseOf{$r\neq\emptyset$.}
        Then $\emptyset\in r$ and $x\in c(\emptyset)$.
        Take a natural number $q$ such that $r=\suc{q}$
            by \obv{natural_zero_or_successor}.
        $R=\suc{\suc{q+q}}$
            by \auto{num_addition_successor_left,num_addition_successor_right}.
        Hence $R\neq\emptyset$ by \auto{suc_neq_emptyset}.
        Thus $\emptyset\in R$
            by \obv{natural_nonzero_contains_zero}.
        $\emptyset\in\suc{r}$
            by \obv{natural_nonzero_contains_zero,omega_successor,suc_neq_emptyset}.
        $d(\emptyset+\emptyset)=c(\emptyset)$
            by \auto{finite_interleave_even}.
        $\emptyset+\emptyset=\emptyset$
            by \obv{omega_empty,num_addition_zero_right}.
        Hence $d(\emptyset)=c(\emptyset)$.
        Thus $x\in d(\emptyset)$.
        Follows by \auto{entry_rank_eq_zero_iff}.
    \end{byCase}
\end{proof}

\begin{lemma}\label{entry_rank_interleave_no_entry_before_inserted}
    Let $r$ be a natural number.
    Suppose $c\in\funs{\suc{r}}{A}$ and $m\in\funs{r}{A}$.
    Let $j\in r$.
    Suppose for every $i\in r$ we have
        $m(i)\subseteq c(\suc{i})$.
    Suppose $\entryRank{c}{r}{x}=\suc{j}$.
    Let $k\in\suc{j+j}$.
    Then $x\notin\apply{\finiteInterleave{c}{m}}{k}$.
\end{lemma}
\begin{proof}
    Let $d=\finiteInterleave{c}{m}$.
    $\dom{c}=\suc{r}$ and $\dom{m}=r$ by \obv{funs_elim}.
    $j$ is a natural number
        by \obv{naturals_is_transitive_set,transitiveset_iff_subseteq,subseteq}.
    $\suc{j}=\entryRank{c}{r}{x}\subseteq r$
        by \obv{entry_rank_subseteq_sentinel}.
    $k\in\finiteEvenIndices{j}\union\finiteOddIndices{j}$
        by \obv{finite_double_index_partition}.
    Hence $k\in\finiteEvenIndices{j}$ or
        $k\in\finiteOddIndices{j}$ by \obv{union_iff}.
    \begin{byCase}
    \caseOf{$k\in\finiteEvenIndices{j}$.}
        Take $i\in\suc{j}$ such that $k=i+i$
            by \obv{finite_even_indices}.
        Hence $i\in\entryRank{c}{r}{x}$.
        Thus $x\notin c(i)$
            by \obv{entry_rank_no_earlier_entry}.
        $i\in r$ by \obv{subseteq}.
        Hence $i\in\suc{r}$ by \auto{suc_iff}.
        $d(i+i)=c(i)$ by \auto{finite_interleave_even}.
    \caseOf{$k\notin\finiteEvenIndices{j}$.}
        Then $k\in\finiteOddIndices{j}$.
        Take $i\in j$ such that $k=\suc{i+i}$
            by \obv{finite_odd_indices}.
        $i\in r$
            by \obv{natural_number_is_ordinal,ordinal_transitivity}.
        $i$ is a natural number
            by \obv{naturals_is_transitive_set,transitiveset_iff_subseteq,subseteq}.
        $i$ and $j$ are ordinals
            by \obv{natural_number_is_ordinal}.
        Hence $\suc{i}\subseteq j$
            by \obv{ordinal_suc_subseteq}.
        $j\subseteq\suc{j}$ by \obv{subseteq_self_suc_intro}.
        Hence $\suc{i}\subseteq\suc{j}$
            by \obv{subseteq_transitive}.
        $\suc{i}$ and $\suc{j}$ are ordinals
            by \obv{suc_ordinal}.
        $\suc{i}\neq\suc{j}$
            by \obv{suc_injective,in_irrefl}.
        Hence $\suc{i}\in\suc{j}$
            by \obv{ordinal_subseteq_implies_elem_or_eq}.
        Thus $x\notin c(\suc{i})$
            by \obv{entry_rank_no_earlier_entry}.
        $m(i)\subseteq c(\suc{i})$.
        Hence $x\notin m(i)$ by \obv{subseteq}.
        $d(\suc{i+i})=m(i)$
            by \auto{finite_interleave_odd}.
    \end{byCase}
\end{proof}

\begin{proposition}\label{entry_rank_interleave_inserted}
    Let $r$ be a natural number.
    Suppose $c\in\funs{\suc{r}}{A}$ and $m\in\funs{r}{A}$.
    Let $j\in r$.
    Suppose for every $i\in r$ we have
        $m(i)\subseteq c(\suc{i})$.
    Suppose $\entryRank{c}{r}{x}=\suc{j}$ and $x\in m(j)$.
    Then
        $\entryRank{\finiteInterleave{c}{m}}{r+r}{x}
            =\suc{j+j}$.
\end{proposition}
\begin{proof}
    Let $d=\finiteInterleave{c}{m}$.
    Let $R=r+r$.
    Let $q=\suc{j+j}$.
    $\dom{c}=\suc{r}$ and $\dom{m}=r$ by \obv{funs_elim}.
    $j$ is a natural number
        by \obv{naturals_is_transitive_set,transitiveset_iff_subseteq,subseteq}.
    $R$ is a natural number by \auto{num_addition_on_naturals}.
    $q\in\finiteOddIndices{r}$
        by \obv{finite_odd_indices}.
    Hence $q\in\finiteEvenIndices{r}\union\finiteOddIndices{r}$
        by \auto{union_iff}.
    Thus $q\in\suc{R}$
        by \auto{finite_double_index_partition}.
    $q\neq R$ by \obv{natural_double_not_successor_double}.
    Hence $q\in R$ by \auto{suc_iff}.
    $d(q)=m(j)$ by \auto{finite_interleave_odd}.
    Hence $x\in d(q)$.
    Thus $q\in\entryRankCandidates{d}{R}{x}$
        by \auto{entry_rank_candidates_iff}.
    For every $i\in q$ we have $x\notin d(i)$
        by \auto{entry_rank_interleave_no_entry_before_inserted}.
    Follows by \auto{entry_rank_eq_from_first_candidate}.
\end{proof}

\begin{proposition}\label{entry_rank_interleave_after_inserted}
    Let $r$ be a natural number.
    Suppose $c\in\funs{\suc{r}}{A}$ and $m\in\funs{r}{A}$.
    Let $j\in r$.
    Suppose for every $i\in r$ we have
        $m(i)\subseteq c(\suc{i})$.
    Suppose $\entryRank{c}{r}{x}=\suc{j}$ and $x\notin m(j)$.
    Then
        $\entryRank{\finiteInterleave{c}{m}}{r+r}{x}
            =\suc{\suc{j+j}}$.
\end{proposition}
\begin{proof}
    Let $d=\finiteInterleave{c}{m}$.
    Let $R=r+r$.
    Let $p=\suc{j+j}$.
    Let $q=\suc{p}$.
    $\dom{c}=\suc{r}$ and $\dom{m}=r$ by \obv{funs_elim}.
    $j$ is a natural number
        by \obv{naturals_is_transitive_set,transitiveset_iff_subseteq,subseteq}.
    Hence $\suc{j}$ is a natural number by \obv{omega_successor}.
    $R$ is a natural number by \auto{num_addition_on_naturals}.
    $j$ and $r$ are ordinals by \obv{natural_number_is_ordinal}.
    Hence $\suc{j}\subseteq r$ by \obv{ordinal_suc_subseteq}.
    $\suc{j}$ and $r$ are ordinals
        by \obv{suc_ordinal,natural_number_is_ordinal}.
    Hence $\suc{j}\in r$ or $\suc{j}=r$
        by \obv{ordinal_subseteq_implies_elem_or_eq}.
    Thus $\suc{j}\in\suc{r}$ by \auto{suc_iff}.
    $q=\suc{j}+\suc{j}$
        by \auto{num_addition_successor_left,num_addition_successor_right}.
    Hence $q\in\finiteEvenIndices{r}$
        by \obv{finite_even_indices}.
    Thus $q\in\finiteEvenIndices{r}\union\finiteOddIndices{r}$
        by \auto{union_iff}.
    Hence $q\in\suc{R}$
        by \auto{finite_double_index_partition}.
    $d(q)=c(\suc{j})$ by \auto{finite_interleave_even}.
    We have $\suc{j}=r$ or
        $\suc{j}\in r$ and $x\in c(\suc{j})$
        by \obv{entry_rank_successor_candidate}.
    Show that $q=R$ or $q\in R$ and $x\in d(q)$.
    \begin{subproof}
        \begin{byCase}
        \caseOf{$\suc{j}=r$.}
            Then $q=R$.
        \caseOf{$\suc{j}\neq r$.}
            Then $\suc{j}\in r$ and $x\in c(\suc{j})$.
            $q\neq R$ by \obv{natural_double_injective}.
            Hence $q\in R$ by \auto{suc_iff}.
            $x\in d(q)$.
        \end{byCase}
    \end{subproof}
    Show that for every $k\in q$ we have $x\notin d(k)$.
    \begin{subproof}
        Fix $k\in q$.
        $k\in p$ or $k=p$ by \auto{suc_iff}.
        \begin{byCase}
        \caseOf{$k\in p$.}
            Follows by
                \auto{entry_rank_interleave_no_entry_before_inserted}.
        \caseOf{$k\notin p$.}
            Then $k=p$.
            $d(p)=m(j)$ by \auto{finite_interleave_odd}.
        \end{byCase}
    \end{subproof}
    Follows by \auto{entry_rank_eq_from_first_candidate}.
\end{proof}

\begin{corollary}\label{entry_rank_interleave_outside_last}
    Let $r$ be a natural number.
    Suppose $c\in\funs{\suc{r}}{A}$ and $m\in\funs{r}{A}$.
    Let $j\in r$.
    Suppose for every $i\in r$ we have
        $m(i)\subseteq c(\suc{i})$.
    Suppose $\entryRank{c}{r}{x}=\suc{j}$ and $\suc{j}=r$.
    Suppose $x\notin c(r)$.
    Then
        $\entryRank{\finiteInterleave{c}{m}}{r+r}{x}=r+r$.
\end{corollary}
\begin{proof}
    $j$ is a natural number
        by \obv{naturals_is_transitive_set,transitiveset_iff_subseteq,subseteq}.
    $m(j)\subseteq c(\suc{j})=c(r)$.
    Hence $x\notin m(j)$ by \obv{subseteq}.
    $\entryRank{\finiteInterleave{c}{m}}{r+r}{x}
        =\suc{\suc{j+j}}$
        by \obv{entry_rank_interleave_after_inserted}.
    $r+r=\suc{\suc{j+j}}$
        by \auto{num_addition_successor_left,num_addition_successor_right}.
\end{proof}

\begin{theorem}\label{entry_rank_under_finite_interleave}
    Let $r$ be a natural number.
    Suppose $c\in\funs{\suc{r}}{A}$ and $m\in\funs{r}{A}$.
    Suppose for every $i\in r$ we have
        $c(i)\subseteq m(i)$ and $m(i)\subseteq c(\suc{i})$.
    Suppose $\entryRank{c}{r}{x}=\suc{j}$.
    Then
        $\entryRank{\finiteInterleave{c}{m}}{r+r}{x}
            =\suc{j+j}$ or
        $\entryRank{\finiteInterleave{c}{m}}{r+r}{x}
            =\suc{\suc{j+j}}$.
\end{theorem}
\begin{proof}
    $j\in r$ by \obv{entry_rank_successor_predecessor}.
    \begin{byCase}
    \caseOf{$x\in m(j)$.}
        Follows by \auto{entry_rank_interleave_inserted}.
    \caseOf{$x\notin m(j)$.}
        Follows by \auto{entry_rank_interleave_after_inserted}.
    \end{byCase}
\end{proof}
